\subsection{HopsFS}
\label{sec:hopsfs}
HopsFS \cite{DBLP:conf/fast/NiaziIHDGR17} is a fork of Apache HDFS created with the explicit goal of solving the biggest scalability and availability limits that are intrinsic to the single-namenode nature of the system: 
\begin{inparaenum}[i)]
    \item amount of metadata limited by the main memory of the machine running the namenode process
    \item number and speed of processors in the machine
    \item amount and latency of bandwidth between the namenode and its clients
    \item coarse grained locking that requires a global lock to alter any piece of metadata
    \item long garbage collection pauses can block the entire process for long periods of time as heap grows
\end{inparaenum}
To do so, HopsFS decouples the responsibility of managing metadata from the namenode and places it in a separate distributed system called MySQL Cluster.
MySQL Cluster is a distributed, consistent (CP), in-memory relational database management system (RDMS) that can be operated and scaled independently from the hadoop cluster(s) it stores metadata for.
Data stored in MySQL Cluster's distributed storage engine (NDB), is divided between nodes participating in the cluster, allowing capacity to be increased by adding more machines to the cluster.
Unlike more traditional RDBMS, where data is stored on disk and only loaded in memory at query time, data in MySQL Cluster is stored in-memory and persisted to disk as a recovery mechanism, allowing very fast query execution.
By moving metadata to such a system, all of the issues regarding the memory limitations of a single system are automatically solved.
The gains are even more significant with regards to the amount of queries per seconds that the system can manage.
Decoupling metadata management from the namenode makes it a stateless component, which can be horizontally scaled and enables downtime-free failover, which is described in the following section.
Furthermore, compared with the approach of having a global lock for all metadata, a relational system \footnote{describe relational database, concept of row and lock on a row} such as MySQL Cluster can have much more fine grained locks allowing, for instance, parallel modification of the information of any number of different files.
Unlike memory-managed applications, MySQL Cluster also does not suffer from garbage collection pauses, avoiding the pitfall in performance as the amount of managed metadata grows larger.

\subsubsection{Multi-namenode architecture}
The namenode, which is now a client of the metadata storage system, performs metadata queries, both in terms of reading and modifying, using an interface called the (meta)Data Access Layer or DAL, which internally connects to the distributed storage system in an efficient fashion.
This allows multiple HopsFS namenodes to run in parallel, each serving a subset of the client requests to the overall system.
The architecture of the resulting system is shown in figure \ref{fig:hopsfs-architecture}.

While most client operations can be directed to any one namenode, the block reports from datanodes and the daemon threads must be handled carefully.
In a Apache HDFS namenode, background daemon threads are responsible for a variety of functions including block re-balancing and lease recovery.
If allowed to run on every node in the cluster, these daemon threads would interfere with each other's work, causing unpredictable results.
To maintain the behaviour of Apache HDFS, HopsFS must guarantee that only one node in the cluster executes daemon threads at any given time.
To this end, all namenodes active in a cluster participate in a leader election algorithm, explained in detail in section \ref{sec:leader-election}, that will elect a single namenode to be leader at any given time.
Additionally, the leader election procedure serves as \emph{failure detector} within the set of namenodes, allowing clients to retrieve a list of live nodes.
Datanodes fetch an updated list of active namenodes before every heartbeat messages, allowing them to keep an updated view of participants in the cluster and the current leader node's identity.
As metadata about the state of datanodes, the action they should execute, and the blocks they contain is stored in MySQL Cluster, any namenode can process heartbeats or block reports from any datanode.
Heartbeats are sent by the datanodes to different namenodes in a round-robin fashion, while the namenode to send block reports to is obtained by performing a RPC call to the leader.
The leader will assign the block report to a namenode that has the capacity to process it and has the lowest number of block reports to process at the moment where the RPC is sent.

\begin{figure}[h]
\caption{Architecture diagram of HopsFS. The metadata storage engine is provided by MySQL Cluster}
\label{fig:hopsfs-architecture}
\centering
\includegraphics[width=0.9\textwidth]{images/placeholder.png}
\end{figure}

\subsubsection{Group membership service and leader election}
\label{sec:leader-election}
As previously mentioned, HopsFS maintains a list of active namenodes in the system and elects one leader which executes functions critical to the proper functioning of the overall system.
Apache HDFS, when configured in High Availability mode (HA) also require such a system and leverages Zookeeper to perform these functions.
HopsFS eliminates the reliance on Zookeeper by leveraging the metadata storage system as shared memory and implementing reliable failure detection and leader election \cite{DBLP:conf/dais/NiaziIBD15} using MySQL Cluster distributed engine (NDB).
\cite{DBLP:journals/jacm/ChandraHT96}

\subsubsection{The (meta)Data Access Layer}
The Data Access Layer, DAL for short, is the Hops component that allows the delegation of metadata handling to the MySQL Cluster database.
To achive this, the component provides two distinct pieces of functionality:
\begin{inparaenum}[i)]
    \item management of the life-cycle of database connection, including various optimizations to reduce network round trips and,
    \item abstractions that allow engineers working on Hops to convert all memory metadata accesses in database operations in a convenient way.
\end{inparaenum}
Formally, the data access layer provides only the interfaces that Hops itself uses to describe accesses to metadata, delegating the implementation of database access to a further library that provides concrete implementations.
Given that in Hops only one such implementation exists (\textit{hops-metadata-dal-impl-ndb}), this chapter will consider both DAL and DAL-implmenetation as a whole without making the distinction explicit.

\paragraph{Connection management}
The DAL provides management of the life-cycle of database connections to a MySQL Cluster cluster.
Specifically, upon configuration, the DAL creates two persistent connectors to the same MySQL Cluster cluster: one that connects to NDB using the native protocol and the ClusterJ java library and one that connects to SQL nodes using the standard JDBC MySQL driver.
The reason to use both a SQL driver and the native NDB protocol is that, while the NDB protocol is very fast at performing primary-key based operations, more complex operations such as joins and deletes are not supported and can only be executed through the SQL nodes.
Given that the performance of Hops is determined mostly by how fast it accesses metadata, the DAL must be as performant as possible.
To achieve better performance, this part of the DAL library implements optimizations aimed at reducing connection overhead, thus allowing a greater number of operations per second.
The main technique for this is connection pooling, which associates each open connection to a thread that will use it for all operations.
By allowing a thread to re-use the same connection for all operations, the overhead of opening the connection is effectively eliminated.
Connections are only closed in case of shutdown of Hops or errors on the connection itself, in which case the connection is re-opened at the next use. 
The connector itself is provided to clients as a global object, accessible to any component that requires it and it is initialized and configured in \texttt{Namenode.java}.

\paragraph{Database access}
Aside from managing database connections, the DAL provides abstractions that are used to convert all memory metadata accesses into accesses to the metadata storage layer.
The main abstraction provided is the \emph{request handler}, a structure that provides information on the type of operation being performed (the \texttt{OPCODE}) and the procedure to execute on the metadata, whether read-only or a modification.
When the handler is executed it performs the procedure in the context of a database transactions where errors will be handled by rolling-back the transaction itself, guaranteeing atomicity of metadata modifications.

The DAL provides two types of request handlers, the \emph{lightweight request handlers} which execute the operations as described above and the \emph{transactional request handlers} that apply most metadata modifications in memory before committing them to the database with the goal of reducing database round-trips.

In a \emph{lightweight request handler}, shown in Listing \ref{lst:lrh}, every modification to the metadata is concretely executed as a database query, causing a large number of network operations.
In case of transaction handlers with very large number of modifications, the network round-trips rapidly become the performance bottleneck.

To increase performance in modification heavy handlers, \emph{transactional request handlers}, shown in Listing \ref{lst:trh}, operate in a different way with the goal of reducing network operations to a minimum.
Transactional request handlers introduce a lock acquisition phase which is executed before the code for the transaction itself.
In this phase, the DAL acquires locks on all the specified rows and materializes them as objects in the DAL memory.
Upon execution, the handler operates on the in-memory representation of the objects either by modifying or deleting existing ones or by creating new ones through the \texttt{EntityManager} class.
At the end of the perform phase, the objects are divided into four categories:
\begin{inparaenum}[1)]
    \item unmodified,
    \item created,
    \item deleted,
    \item modified
\end{inparaenum}
and the required operations are executed in batch on the database.
The perform phase is, therefore, still executed in the context of a database transaction, with the possibility of rollback in case of errors, but all operations on the database are executed at the end.
Given that all of the materialized rows are locked for the duration of the database, there can be no conflicts upon commit at the end of the handler.
Note that the handler can request read locks as well as write locks and, in that case, the rows locked in read mode cannot be modified.

Replacing all memory accesses with transaction handlers which acquire the minimum amount of locks required to perform the operation, HopsFS achieves a much more granular level of concurrency compared to the in-memory global lock, which allows it to execute a much greater number of concurrent operations.


\begin{lstlisting}[label={lst:trh}, caption={Transactional request handler for the rename operation}]
OperationType opType
if(isUsingSubTreeLocks) {
    opType = HDFSOperationType.SUBTREE_RENAME;
} else {
    opType = HDFSOperationType.RENAME;
}
new HopsTransactionalRequestHandler(opType, src) {
    @Override
    public void acquireLock(TransactionLocks locks)
    throws IOException {
        LockFactory lf = LockFactory.getInstance();
        locks.add(lf.getRenameINodeLock(
            nameNode, INodeLockType.WRITE_ON_TARGET_AND_PARENT,
            INodeResolveType.PATH, true, src, dst))
        .add(lf.getBlockLock())
        .add(lf.getBlockRelated(
            BLK.RE, BLK.CR, BLK.UC,
            BLK.UR, BLK.IV, BLK.PE, BLK.ER));
        if (dir.isQuotaEnabled()) {
            locks.add(lf.getQuotaUpdateLock(
                true, src, dst));
        }
        if (!isUsingSubTreeLocks) {
            locks.add(lf.getLeaseLock(
                LockType.WRITE))
            .add(lf.getLeasePathLock(
                LockType.READ_COMMITTED));
        } else {
            locks.add(lf.getLeaseLock(
                LockType.WRITE))
            .add(lf.getLeasePathLock(
                LockType.WRITE, src));
        }
        if (erasureCodingEnabled) {
            locks.add(lf.getEncodingStatusLock(
                LockType.WRITE, dst));
        }
    }

    @Override
    public Object performTask(StorageConnector connector)
    throws IOException {
        if (NameNode.stateChangeLog.isDebugEnabled()) {
            NameNode.stateChangeLog.debug(
                "DIR* NameSystem.renameTo: with options - " + src + " to " + dst);
        }

        if (isInSafeMode()) {
            throw new SafeModeException("Cannot rename " + src, safeMode);
        }
        if (!DFSUtil.isValidName(dst)) {
            throw new InvalidPathException("Invalid name: " + dst);
        }
        for (MetadataLogEntry logEntry: logEntries) {
            EntityManager.add(logEntry);
        }

        for (Options.Rename op: options) {
            if (op == Rename.KEEP_ENCODING_STATUS) {
                INode[] srcNodes = dir.getRootDir()
                    .getExistingPathINodes(src, false);
                INode[] dstNodes = dir.getRootDir()
                    .getExistingPathINodes(dst, false);
                INode srcNode = 
                    srcNodes[srcNodes.length - 1];
                INode dstNode = 
                    dstNodes[dstNodes.length - 1];
                EncodingStatus status = EntityManager.find(
                    EncodingStatus.Finder.ByInodeId, dstNode.getId());
                EncodingStatus newStatus = new EncodingStatus(status);
                newStatus.setInodeId(srcNode.getId());
                EntityManager.add(newStatus);
                EntityManager.remove(status);
                break;
            }
        }

        removeSubTreeLocksForRenameInternal(
            src, isUsingSubTreeLocks, subTreeLockDst);

        dir.renameTo(
            connector, src, dst, srcNsCount,
            srcDsCount, dstNsCount, dstDsCount, options);
        return null;
    }
}.handle(this);
\end{lstlisting}

\begin{lstlisting}[label={lst:lrh}, caption={Lightweight request handler }]
new LightWeightRequestHandler(UsersOperationsType.GET_USER_GROUPS) {
    @Override
    public Object performTask(StorageConnector connector)
    throws IOException {
        boolean transactionActive = connector.isTransactionActive();

        if (!transactionActive) {
            connector.beginTransaction();
        }

        Integer userId = cache.getUserId(userName);
        
        User user;
        if(userId == null) {
            user = userDataAccess.getUser(userName);
        } else {
            user = userDataAccess.getUser(userId);
        }

        if (user == null) {
            return null;
        }

        List<Group> groups = userGroupDataAccess
            .getGroupsForUser(user.getId());

        if (!transactionActive) {
            connector.commit();
        }

        return new Pair<User, List<Group>>(user, groups);
}.handle(this);
\end{lstlisting}